% ============================================================
\section{Cross-Domain Validation: Forecasting}
\label{sec:forecasting}
% ============================================================

The IoT anomaly detection results raise a fundamental question:
does fine-tuning destroy \emph{genuinely useful} pre-trained features,
or were those features never useful to begin with?
In IoT anomaly detection, a logistic regression baseline on 72-dimensional
hand-crafted features achieves AUC$=$0.630, statistically indistinguishable
from the best Moirai condition (C: 0.629$\pm$0.022).
This parity makes it difficult to disentangle catastrophic forgetting
from a low-signal task ceiling.

To resolve this ambiguity, we replicate the forgetting diagnosis on
\textbf{ETTh2 time-series forecasting}~\cite{zhou2021informer},
a domain where Moirai's pre-trained features are demonstrably valuable:
zero-shot Moirai outperforms a Linear baseline by 28\% at $h{=}96$
and 45\% at $h{=}192$ (Section~\ref{sec:forecasting:gate}).
If forgetting replicates here, where pre-trained features provide
measurable advantage, the mechanism is general, not an artefact
of low task signal.

% ============================================================
\subsection{Domain Selection and Pre-Training Value Gate}
\label{sec:forecasting:gate}
% ============================================================

We require a forecasting benchmark where Moirai zero-shot performance
substantially exceeds simple baselines, establishing that pre-trained
features provide genuine value.
Following \citet{woo2024moirai} Table~6, we evaluate on ETTh2
(Electricity Transformer Temperature, Hourly, 7 features, 17{,}420
timesteps, standard 12/4/4 month train/val/test split).

\textbf{Gate criterion:} Moirai zero-shot MSE must improve over
Linear regression by $>$20\%.

\textbf{Results.}
Using the published Moirai-Small (13.8M parameters) with median point
forecasts (robust to heavy-tailed mixture distribution components;
see Appendix~\ref{app:moirai_median}):

\begin{center}
\small
\begin{tabular}{lccc}
\toprule
& Linear MSE & Moirai MSE & Gap \\
\midrule
$h=96$  & 0.373 & \textbf{0.269} & $+$28.0\% \\
$h=192$ & 0.544 & \textbf{0.299} & $+$45.0\% \\
\bottomrule
\end{tabular}
\end{center}

Both horizons pass the 20\% gate.
Pre-trained features are demonstrably useful on this task, unlike
IoT anomaly detection, where the frozen encoder barely exceeds
zero-shot (AUC: 0.549 vs.\ 0.481).

% ============================================================
\subsection{Experimental Setup}
\label{sec:forecasting:setup}
% ============================================================

We mirror the IoT forgetting diagnosis (Section~\ref{sec:experiments})
with three conditions:

\begin{itemize}
\item \textbf{B (NLL-only):} Fine-tune all 13.8M parameters using
  negative log-likelihood on context$+$target sequences (patch size 32).
\item \textbf{C (NLL$+$SupCon):} Add temporal supervised contrastive
  loss ($\lambda{=}0.5$) with 8 temporal clusters as pseudo-labels.
  Sequences within the same temporal regime are positives; different
  regimes are negatives.
\item \textbf{D (Frozen encoder):} Freeze the Moirai encoder;
  train only the output head.
\end{itemize}

\textbf{Training:} 500 samples, 20 epochs, AdamW (lr$=$1e-4),
cosine annealing, gradient clipping (1.0).
\textbf{Evaluation:} Median of 20 forecast samples on held-out
validation set (300 sequences), normalised using training statistics.
\textbf{Diagnosis:} CKA between pre-trained and fine-tuned encoder
representations, weight drift ($\ell_2$ from initialisation).
All results report mean$\pm$std over 5 seeds
(42, 123, 456, 789, 1234).

% ============================================================
\subsection{Forgetting Replicates on High-Signal Task}
\label{sec:forecasting:results}
% ============================================================

\begin{table}[t]
\centering
\caption{Representation drift diagnosis on ETTh2 forecasting (mean $\pm$ std across 10 seeds).
Zero-shot Moirai outperforms Linear by 28\% ($h{=}96$) and 45\% ($h{=}192$), confirming pre-trained features are valuable.
Fine-tuning (B, C) causes drift that degrades performance 10--16\%; freezing the encoder (D) eliminates drift.}
\label{tab:forecasting_forgetting}
\small
\begin{tabular}{l cccc cccc}
\toprule
& \multicolumn{4}{c}{$h=96$ (ZS MSE: 0.126)} & \multicolumn{4}{c}{$h=192$ (ZS MSE: 0.158)} \\
\cmidrule(lr){2-5} \cmidrule(lr){6-9}
Condition & MSE & Forg.\% & CKA & Drift & MSE & Forg.\% & CKA & Drift \\
\midrule
B: NLL-only    & .140{\tiny$\pm$.014} & +11.1  & .937{\tiny$\pm$.028} & 3.96{\tiny$\pm$.56}
               & .183{\tiny$\pm$.020} & +15.7 & .970{\tiny$\pm$.009} & 4.12{\tiny$\pm$.32} \\
C: NLL+SupCon  & .138{\tiny$\pm$.010} & +10.1  & .949{\tiny$\pm$.014} & 3.72{\tiny$\pm$.43}
               & .178{\tiny$\pm$.011} & +12.9 & .964{\tiny$\pm$.015} & 3.97{\tiny$\pm$.34} \\
D: Frozen enc. & \textbf{.126}{\tiny$\pm$.006} & \textbf{+0.5}  & \textbf{1.00}{\tiny$\pm$.000} & \textbf{1.11}{\tiny$\pm$.15}
               & \textbf{.155}{\tiny$\pm$.002} & \textbf{$-$1.7} & \textbf{1.00}{\tiny$\pm$.000} & \textbf{1.30}{\tiny$\pm$.18} \\
\bottomrule
\end{tabular}
\end{table}


Table~\ref{tab:forecasting_forgetting} shows the central result:
\textbf{catastrophic forgetting replicates on ETTh2}, where pre-trained
features are demonstrably valuable.

\textbf{Unfrozen conditions degrade.}
Both B (NLL-only) and C (NLL$+$SupCon) show significant forgetting:
$+$8.9\% and $+$9.0\% at $h{=}96$, increasing to $+$12.6\% and
$+$13.1\% at $h{=}192$.
The NLL training loss decreases monotonically throughout, confirming
that the model is learning, but what it learns destroys the pre-trained
representations that enable accurate zero-shot forecasting.

\textbf{Frozen encoder preserves performance.}
Condition D maintains near-zero forgetting ($+$0.9\% at $h{=}96$,
$-$1.7\% at $h{=}192$) with perfect CKA ($=$1.000) and $4\times$
lower weight drift.
This confirms the mechanism: fine-tuning overwrites pre-trained features;
freezing prevents this.

\textbf{Forgetting scales with horizon.}
At $h{=}192$ (where Moirai's pre-training advantage is 45\%), forgetting
is more severe (12--13\%) than at $h{=}96$ (28\% advantage, 9\% forgetting).
Longer horizons require the model to leverage more of its pre-trained
temporal knowledge, making it more vulnerable to catastrophic forgetting.

\textbf{Contrastive loss does not prevent forgetting.}
Unlike the common hypothesis that contrastive learning regularises
representations, condition C shows comparable or slightly worse
forgetting than B.
The temporal contrastive objective pushes representations toward
task-specific clustering, which does not preserve the distributional
structure that makes zero-shot forecasting effective.

% ============================================================
\subsection{Cross-Domain Comparison}
\label{sec:forecasting:comparison}
% ============================================================

Comparing the forgetting diagnosis across domains:

\begin{center}
\small
\begin{tabular}{lcc}
\toprule
& IoT Anomaly Det. & ETTh2 Forecasting \\
\midrule
Pre-training value        & Marginal (0.549 vs.\ 0.481) & Strong (28--45\%) \\
Forgetting (unfrozen)     & $+$7--15\% AUC drop & $+$9--13\% MSE rise \\
Frozen encoder recovery   & Yes (AUC: 0.549--0.606) & Yes (MSE: $\approx$zero-shot) \\
CKA degradation           & 0.426$\to$0.183 & 0.939--0.971 \\
LogReg parity             & Yes (0.630 $\approx$ 0.629) & N/A (no features) \\
\bottomrule
\end{tabular}
\end{center}

The key insight: \textbf{forgetting occurs regardless of whether
pre-trained features are useful.}
In IoT, the marginal pre-trained signal (AUC: 0.549 vs.\ 0.481)
is destroyed, but this is difficult to distinguish from a task ceiling
since LogReg matches the fine-tuned model.
In forecasting, pre-trained features provide a large, measurable
advantage (28--45\% MSE improvement) that fine-tuning systematically
degrades.

This cross-domain replication rules out the alternative interpretation
that IoT forgetting was merely ``the task has low signal and models
overfit the synthetic distribution.''
The mechanism is general: extended fine-tuning with NLL (and optionally
contrastive) loss overwrites pre-trained representations, and
freezing the encoder is an effective mitigation across both domains.
